\documentclass{article}
\usepackage{amssymb}
\usepackage{amsmath}
\usepackage{amsthm}
\usepackage{tikz-cd}
\usepackage{mathpartir}
\usepackage{stmaryrd}
\usepackage{parskip}
\usepackage{xcolor}
\usepackage{array}
\usepackage[margin=1.0in]{geometry}
\usepackage{comment}

\definecolor{myorange}{RGB}{242, 114, 0}
\definecolor{myblue}{RGB}{0,162,255}

\newcommand{\semptb}[1]{\morbisimid{#1}}
\newcommand{\ltsqdelay}[2]{\mathrel{{\ltdyn}\text{ptb}^{#1}_{#2}}}

% restriction of a relation R along a vertical morphism f on the left
% \newcommand{\restr}[2]{\mathrel{{#2}^*{#1}}}
% \newcommand{\restrl}[2]{\restr{#1}{#2}^l} 
% \newcommand{\restrr}[2]{\restr{#1}{#2}^r}

\newcommand{\restrl}[2]{\mathrel{{#2}^*{#1}}}
\newcommand{\restrr}[2]{\mathrel{{#1}^*{#2}}}

%%%%%%%%%%%%%%%%%%%%%%%%
% theorem environments
%%%%%%%%%%%%%%%%%%%%%%%%
\newtheorem{theorem}{Theorem}[section]
\newtheorem{lemma}{Lemma}[section]
\newtheorem{nonnum-theorem}{Theorem}[section]
\newtheorem{corollary}{Corollary}[section]

\theoremstyle{definition}
\newtheorem{definition}{Definition}[section]

\theoremstyle{remark}
\newtheorem{remark}{Remark}
%%%%%%%%%%%%%%%%%%%%%%%%%%%%%%%%%%%%%%%%%%%%%%%%%%%

\newcommand{\To}{\Rightarrow}
\newcommand{\inl}{\mathsf{inl}}
\newcommand{\inr}{\mathsf{inr}}
\newcommand{\alt}{\mathrel{\bf \,\mid\,}}

\newcommand{\ob}{\text{Ob}}


\newcommand{\extlc}{\text{Ext-}\lambda}
\newcommand{\extlcm}{\text{Ext-}\lambda^{-\text{trans}}}
\newcommand{\extlcmm}{\text{Ext-}\lambda^{-\text{trans}-\text{cast}}}
\newcommand{\extlcprime}{\text{Ext-}\lambda'}
\newcommand{\intlc}{\text{Int-}\lambda}
\newcommand{\intlcbisim}{\text{Int$_\approx$-}\lambda}
\newcommand{\erase}[1]{\lfloor {#1} \rfloor}


\newcommand{\uarrowl}{\mathrel{\rotatebox[origin=c]{-30}{$\leftarrowtail$}}}
\newcommand{\uarrowr}{\mathrel{\rotatebox[origin=c]{60}{$\leftarrowtail$}}}
\newcommand{\darrowl}{\mathrel{\rotatebox[origin=c]{30}{$\twoheadleftarrow$}}}
\newcommand{\darrowr}{\mathrel{\rotatebox[origin=c]{120}{$\twoheadleftarrow$}}}
\newcommand{\vuarrow}{\mathrel{\rotatebox[origin=c]{-90}{$\leftarrowtail$}}}
\newcommand{\vdarrow}{\mathrel{\rotatebox[origin=c]{90}{$\twoheadleftarrow$}}}

% Types, terms, and precision 

\newcommand{\dyn}{?}
\newcommand{\nat}{\text{Nat}}
\newcommand{\bool}{\text{Bool}}
\newcommand{\ra}{\rightharpoonup}
\newcommand{\Ret}[1]{\mathsf{Ret\,}{#1}}
\newcommand{\hole}[1]{\bullet \colon {#1}}
\newcommand{\dyntodyn}{\dyn \ra\, \dyn}


\newcommand{\up}[2]{\langle{#2}\uarrowl{#1}\rangle}
\newcommand{\dn}[2]{\langle{#1}\darrowl{#2}\rangle}

\newcommand{\upc}[1]{\text{up}\,{#1}\,}
\newcommand{\dnc}[1]{\text{dn}\,{#1}\,}


% \newcommand{\ret}{\mathsf{ret}}
\newcommand{\err}{\mho}
\newcommand{\zro}{\textsf{zro}}
\newcommand{\suc}{\textsf{suc}}
\newcommand{\lda}[2]{\lambda {#1} . {#2}}

\newcommand{\injarr}[1]{\textsf{Inj}_\ra ({#1})}
\newcommand{\injnat}[1]{\textsf{Inj}_\text{nat} ({#1})}
\newcommand{\casenat}[4]{\text{Case}_\text{nat} ({#1}) \{ \text{no} \to {#2} \alt \text{nat}({#3}) \to {#4} \}}
\newcommand{\casearr}[4]{\text{Case}_\ra ({#1}) \{ \text{no} \to {#2} \alt \text{fun}({#3}) \to {#4} \}}
\newcommand{\casedyn}[5]{\text{Case}_\Dyn ({#1}) \{ \text{nat}({#2}) \to {#3} \alt \text{fun}({#4}) \to {#5} \}}
%\newcommand{\bind}[3]{\text{var } {#1} = {#2} \text{ in } {#3}}
\newcommand{\matchnat}[4]{\text{match } {#1} \text{ with } \{ \textsf {zero} \Rightarrow {#2} \alt \textsf{ suc } {#3} \Rightarrow {#4} \}}

\newcommand{\refl}{\text{refl}}

\newcommand{\Lift}{\text{Lift}}

%\newcommand{\rel}{\circ\hspace{-4px}-\hspace{-4px}\bullet}
\newcommand{\rel}{\mathrel{\circ\mkern-6mu-\mkern-6mu\bullet}}
% \newcommand{\wand}{\mathrel{-\mkern-6mu*}}
\newcommand{\wand}{\mathrel{\multimap}}
\newcommand{\ltdyn}{\sqsubseteq}
\newcommand{\gtdyn}{\sqsupseteq}
\newcommand{\equidyn}{\mathrel{\gtdyn\ltdyn}}
\newcommand{\gamlt}{\Gamma^\ltdyn}
\newcommand{\deltalt}{\Delta^\ltdyn}
\newcommand{\relcomp}{\odot}

\newcommand{\hasty}[3]{{#1} \vdash {#2} \colon {#3}}
\newcommand{\vhasty}[3]{{#1} \vdash^v {#2} \colon {#3}}
\newcommand{\phasty}[3]{{#1} \vdash^c {#2} \colon {#3}}
\newcommand{\etmprec}[4]{{#1} \vdash {#2} \ltdyn_e {#3} \colon {#4}}
\newcommand{\itmprec}[4]{{#1} \vdash {#2} \ltdyn_i {#3} \colon {#4}}
\newcommand{\etmequidyn}[4]{{#1} \vdash {#2} \equidyn_e {#3} \colon {#4}}
\newcommand{\itmequidyn}[4]{{#1} \vdash {#2} \equidyn_i {#3} \colon {#4}}
\newcommand{\synbisim}{\approx_{\text{syn}}}

\newcommand{\Dwn}{\Downarrow}
\newcommand{\qte}[1]{\text{quote}({#1})}

\newcommand{\elab}[1]{\text{Elab}({#1})}

% Perturbations
\newcommand{\pertp}{\text{Pert}^\text{P}}
\newcommand{\perte}{\text{Pert}^\text{E}}
\newcommand{\pertdyn}[2]{\text{pert-dyn}({#1}, {#2})}
\newcommand{\delaypert}[1]{\text{delay-pert}({#1})}

\newcommand{\pertc}{\text{Pert}_{\text{C}}}
\newcommand{\pertv}{\text{Pert}_{\text{V}}}


% SGDT and Intensional Stuff

\newcommand{\later}{{\vartriangleright}}
\newcommand{\laterhs}{{\later}}
\newcommand{\type}{\texttt{Type}}
\newcommand{\lob}{\text{L\"{o}b}}
\newcommand{\tick}{\mathsf{tick}}
\newcommand{\nxt}{\mathsf{next}}
\newcommand{\fix}{\mathsf{fix}}
\newcommand{\kpa}{\kappa}

% Model-related stuff
\newcommand{\calV}{\mathcal{V}}
\newcommand{\calE}{\mathcal{E}}
\newcommand{\calVk}{\mathcal{V}_k}
\newcommand{\calEk}{\mathcal{E}_k}
\newcommand{\tok}{\mathrel{\mathop{\to}\limits^{\textrm{\footnotesize k}}}}
\newcommand{\timesk}{\mathrel{\mathop{\times}\limits^{\textrm{k}}}}
\newcommand{\Fk}{\mathrel{\mathop{F}\limits^{\textrm{k}}}}
\newcommand{\Uk}{\mathrel{\mathop{U}\limits^{\textrm{k}}}}

\newcommand{\op}[1]{{#1}^{\textrm{op}}}
\newcommand{\calC}{\mathcal{C}}
\newcommand{\Set}{\mathsf{Set}}
\newcommand{\ErrDom}{\mathsf{ErrDom}}
\newcommand{\Yo}{\mathsf{Yo}}
\newcommand{\Hom}{\mathsf{Hom}}
\newcommand{\calS}{\mathcal{S}}
\newcommand{\gfix}{\texttt{gfix}}
\newcommand{\qfix}{\texttt{qfix}}
\newcommand{\calU}{\mathcal{U}}
\newcommand{\laterhat}{\widehat{\later}}
\newcommand{\El}{\mathsf{El}}
\newcommand{\Clock}{\mathsf{Clock}}

\newcommand{\Machine}[1]{\mathsf{Machine}\, {#1}}


% Predomains and EP pairs
\newcommand{\Nat}{\mathsf{Nat}}
\newcommand{\Dyn}{\mathsf{Dyn}}
\newcommand{\ty}[1]{\langle {#1} \rangle}
\newcommand{\li}{L_\mho}
\newcommand{\liclk}[1]{L_\mho [{#1}]}

\newcommand{\ext}[2]{\text{ext}\,{#1}\,{#2}}
\newcommand{\map}[2]{\text{map}\,{#1}\,{#2}}

\newcommand{\ltls}{\ltdyn}
\newcommand{\bisim}{\approx}
\newcommand{\semlt}{\le}
\newcommand{\semltbad}{\lesssim}

%\newcommand{\injarr}{\textsf{Inj}_\to}
%\newcommand{\injnat}{\textsf{Inj}_\mathbb{N}}

\newcommand{\id}{\mathsf{id}}
\newcommand{\ep}{\leadsto}

\newcommand{\emb}[2]{\mathsf{emb}_{#1}({#2})}
\newcommand{\proj}[2]{\mathsf{proj}_{#1}({#2})}

\newcommand{\monto}{\to_m}


% Notation for wait functions
\newcommand{\wre}{w_r^e}
\newcommand{\wle}{w_l^e}
\newcommand{\wrp}{w_r^p}
\newcommand{\wlp}{w_l^p}

\newcommand{\sem}[1]{\llbracket {#1} \rrbracket}
\newcommand{\semgl}[1]{{\sem{#1}}^\text{gl}}


% Denotational model
\newcommand{\errdom}{\textsf{ErrDom}}

\newcommand{\ptb}{\text{ptb}}
\newcommand{\push}{\text{push}}
\newcommand{\pull}{\text{pull}}

\newcommand{\pv}{P^{\mathcal{V}}}
\newcommand{\pe}{P^{\mathcal{E}}}
\newcommand{\ptbv}{\text{ptb}^\mathcal{V}}
\newcommand{\ptbe}{\text{ptb}^\mathcal{E}}

\newcommand{\Ppure}{P}
\newcommand{\Pk}{P^K}
\newcommand{\ptbk}{\text{ptb}^K}


\newcommand{\upf}{\text{up}}
\newcommand{\dnf}{\text{dn}
}
\newcommand{\arr}{\to}
\newcommand{\comp}{\bullet}

\newcommand{\ltsq}[2]{\mathrel{\ltdyn^{#1}_{#2}}}
\newcommand{\ltsqbisim}[2]{\mathrel{{\widetilde{\ltdyn}}^{#1}_{#2}}}
\newcommand{\ltbisim}{\mathrel{\widetilde{\ltdyn}}}
\newcommand{\vf}{\mathcal{V}_f}
\newcommand{\vsq}{\mathcal{V}_{sq}}
\newcommand{\ef}{\mathcal{E}_f}
\newcommand{\esq}{\mathcal{E}_{sq}}

\newcommand{\morbisimid}[1]{\text{Endo}_\bisim{#1}}


\newcommand{\sv}{s_{\mathcal{V}}}
\newcommand{\tv}{t_{\mathcal{V}}}
\newcommand{\rv}{r_{\mathcal{V}}}
\newcommand{\se}{s_{\mathcal{E}}}
\newcommand{\te}{t_{\mathcal{E}}}
\newcommand{\re}{r_{\mathcal{E}}}

\newcommand{\Ff}{F_f}
\newcommand{\Fsq}{F_{sq}}
\newcommand{\Uf}{U_f}
\newcommand{\Usq}{U_{sq}}
\newcommand{\arrf}{\arr_f}
\newcommand{\arrsq}{\arr_{sq}}

\newcommand{\vsim}{\mathcal{V}_{\bisim}}
\newcommand{\esim}{\mathcal{E}_{\bisim}}
\newcommand{\svsim}{s_{\mathcal{V}}^\bisim}
\newcommand{\tvsim}{t_{\mathcal{V}}^\bisim}
\newcommand{\rvsim}{r_{\mathcal{V}}^\bisim}
\newcommand{\sesim}{s_{\mathcal{E}}^\bisim}
\newcommand{\tesim}{t_{\mathcal{E}}^\bisim}
\newcommand{\resim}{r_{\mathcal{E}}^\bisim}



\newcommand{\ve}{\mathcal{V}_e}
\newcommand{\ee}{\mathcal{E}_e}

\newcommand{\vr}{\mathcal{V}_r}
\newcommand{\er}{\mathcal{E}_r}

\newcommand{\relatedin}[3]{{#1}\, {#3}\, {#2}}
\newcommand{\binrel}[1]{\mathbin{#1}}


\newcommand{\da}{\downarrow}

\newcommand{\ret}{\text{ret}}
\newcommand{\bind}[3]{{#1} <- {#2}\, ; \, {#3}}

\newcommand{\piv}{\Pi^\mathcal{V}}
\newcommand{\pie}{\Pi^\mathcal{E}}

\newcommand{\upl}{\textsc{UpL}}
\newcommand{\upr}{\textsc{UpR}}
\newcommand{\dnl}{\textsc{DnL}}
\newcommand{\dnr}{\textsc{DnR}}

\newcommand{\delre}{\delta^{r,e}}
\newcommand{\delle}{\delta^{l,e}}
\newcommand{\delrp}{\delta^{r,p}}
\newcommand{\dellp}{\delta^{l,p}}

\newcommand{\qordeq}{\bisim}

\newcommand{\inat}{\text{Inj}_\mathbb{N}}
\newcommand{\itimes}{\text{Inj}_\times}
\newcommand{\iarr}{\text{Inj}_\to}

\newcommand{\tnat}{\mathsf{nat}}
\newcommand{\ttimes}{\mathsf{times}}
\newcommand{\tfun}{\mathsf{fun}}

\newcommand{\cased}[7]{
    \begin{align*}
    \text{Case}_D ({#1}) &\text{ of } \{ \\
        &\alt \text{nat}({#2}) \to {#3}  \\
        &\alt \text{times}({#4}) \to {#5}  \\
        &\alt \text{fun}({#6}) \to {#7} \}
    \end{align*}
}

\newcommand{\inone}{\mathsf{in}_1}
\newcommand{\intwo}{\mathsf{in}_2}
\newcommand{\inthree}{\mathsf{in}_3}
\newcommand{\infour}{\mathsf{in}_4}
\newcommand{\infive}{\mathsf{in}_5}


\newcommand{\delay}{\text{Delay}}
\newcommand{\ledelay}{\le^{\text{Del}}}
\newcommand{\bisimdelay}{\bisim^{\text{Del}}}
\newcommand{\tnow}{\mathsf{now}}
\newcommand{\tlater}{\mathsf{later}}
\newcommand{\Da}{\Downarrow}


\begin{document}

\title{Revisiting the Notion of Representability}
\author{Eric Giovannini}

\maketitle

Our current notion of quasi-representability of relations involves
perturbations associated to each of our representable relations.  This
leads to two complications with our semantics:

\begin{enumerate}
  
% \item Composition of quasi-representable relations requires the
%   push-pull property. We will revisit this point later.
  
\item The projection associated to $F(cc')$ is not simply the composition of the
  projections for $Fc$ and $Fc'$ (unlike the original term semantics); it
  instead involves an additional perturbation. The additional perturbation is
  inserted in the relational semantics as part of the compositional construction
  of quasi-right-representability of $F(cc')$ given quasi-left-representability
  of $c$ and $c'$.
  %
  For a specific example where this perturbation results in observably
  different behavior compared to the original term semantics, see the Appendix
  of this document.

\item The (quasi-)order equivalence of two quasi-representable relations does
  not follow from sharing the same embedding and projection morphisms. Instead,
  we needed to define a separate primitive notion of quasi-equivalence where we
  have squares with perturbations on the side.
  
\end{enumerate}


There is an alternative, more principled approach to perturbaions and
quasi-representable relations. The idea is to define a double category such that
what were \emph{quasi-representable} relations in the original double category
of value objects/morphisms/relations/squares are now actually
\emph{representable} relations in the usual sense (i.e. with identity morphisms
rather than perturbations).

%% To do this, we will formulate a new double category (actually two, one
%% for the values and one for the computations).

There are two key differences between this formulation of perturbations and the
original. First, in the new double setup, the vertical morphisms come equipped
with an \emph{action} on perturbations. Second, the squares in the new double
category allow for an explicit synchronization of the vertical morphisms by a
perturbation on both sides. We elaborate on these differences below, but first,
we must make an important change to the definition of value and computation
types, requiring that the monoid of syntactic perturbations be commutative. We
will see why this is needed later, but for now note the new definition of
value/computation objects.

\begin{definition}
  A \emph{value object} consists of a predomain $A$, a
  \textbf{commutative} monoid of perturbations $M_A$, and an
  homomorphism of monoids $i_A : M_A \to \semptb{A}$.

  A \emph{computation object} is defined similarly: an error domain $B$,
  a \textbf{commutative monoid} $M_B$ and a homomorphism of monoids $i_B$.
\end{definition}

Note that $\semptb{A}$ is not a commutative monoid, so we will need a
more complicated universal property to define a homomorphism from a
freely-constructed commutative monoid into $\semptb{A}$. We will
discuss this further below.

%%%%%%%%%%%%%%%%%%%%%%%%%%%%%%%%%%%%%%%%%%%%%%%%%%%%%%%%%%%%%%%%%%%%%%%%%%%%%
\subsection{Vertical Morphisms}

The vertical morphisms are defined as follows:

\begin{definition}
  Given value objects $A = (|A|, M_A, i_A)$ and $A' = (|A'|, M_{A'}, i_{A'})$, a
  \emph{perturbation-preserving morphism} consists of
  \begin{itemize}
    \item A morphism of the underlying predomains $f : |A| \to |A'|$
    \item Two homomorphisms of monoids $h^\to : M_A \to M_{A'}$ and
      $h^\leftarrow : M_{A'} \to M_A$
    \item Two coherence conditions: for all $m \in M_A$, we have 
    $f \circ i_A(m) = i_{A'}(h^\to(m)) \circ f$:

      % https://q.uiver.app/#q=WzAsNCxbMCwwLCJBIl0sWzAsMSwiQSJdLFsxLDAsIkEnIl0sWzEsMSwiQSciXSxbMCwyLCJmIl0sWzIsMywiaV97QSd9KGhee1xcdG99KG0pKSJdLFswLDEsImlfQShtKSIsMl0sWzEsMywiZiIsMl1d
      \[\begin{tikzcd}[ampersand replacement=\&]
              A \& {A'} \\
              A \& {A'}
              \arrow["f", from=1-1, to=1-2]
              \arrow["{i_A(m)}"', from=1-1, to=2-1]
              \arrow["{i_{A'}(h^{\to}(m))}", from=1-2, to=2-2]
              \arrow["f"', from=2-1, to=2-2]
      \end{tikzcd}\]

    and likewise for all $m' \in M_{A'}$, we have 
    $f \circ i_A(h^\leftarrow(m')) = i_{A'}(m') \circ f$:

      % https://q.uiver.app/#q=WzAsNCxbMCwwLCJBIl0sWzEsMCwiQSciXSxbMCwxLCJBIl0sWzEsMSwiQSciXSxbMCwxLCJmIl0sWzIsMywiZiIsMl0sWzEsMywiaV97QSd9KG0nKSJdLFswLDIsImlfQShoXntcXGxlZnRhcnJvd30obScpKSIsMl1d
      \[\begin{tikzcd}[ampersand replacement=\&]
              A \& {A'} \\
              A \& {A'}
              \arrow["f", from=1-1, to=1-2]
              \arrow["{i_A(h^{\leftarrow}(m'))}"', from=1-1, to=2-1]
              \arrow["{i_{A'}(m')}", from=1-2, to=2-2]
              \arrow["f"', from=2-1, to=2-2]
      \end{tikzcd}\]
  \end{itemize}

  We make the analogous definition for computation objects.
\end{definition}

\vspace{1ex}

\begin{remark}
Not all predomain/error domain morphisms extend to an action on perturbations in
this way. The key counterexample is the arrow type constructor. The issue is
that the perturbations of $A \to B$ are defined as $M_A^{op} \oplus M_B$, i.e.,
both $M_A$ and $M_B$ are in covariant positions. But the universal property for
the arrow type says that
%
\[ \mathcal{V}(\Gamma, U(A \to B)) \cong \mathcal{V}(\Gamma \times A, UB) \]
%
On the left, the perturbations are monoid homomorphisms $M_\Gamma \multimap
(M_A^{op} \oplus M_B)$ whereas on the right they are monoid homomorphisms
$(M_\Gamma \oplus M_A) \multimap M_B$.
\end{remark}


%% With this new definition of morphism, the next step is to show that
%% the morphisms of predomains/error domains used in the graduality proof
%% extend to perturbations.


%%%%%%%%%%%%%%%%%%%%%%%%%%%%%%%%%%%%%%%%%%%%%%%%%%%%%%%%%%%%%%%%%%%%%%%%%%%%%
\subsection{Squares}

A square in this double category will be defined as a square in the usual sense,
but where we are allowed to synchronize the vertical morpisms on either side
with a perturbation. More concretely:

\begin{definition}\label{def:perturbed-square}
Let $f : A_i \to A_o$ and $g : A_i' \to A_o'$ be
perturbation-preserving predomain morphisms and let $c_i : A_i \rel
A_i'$ and $c_o : A_o \rel A_o'$ be predomain relations.
%
We say that there is a \textbf{perturbed square} between $f$ and $g$, written $f
\ltsqdelay{c_i}{c_o} g$, if there \emph{merely exist} perturbations $m_f$ and
$m_g$ such that $i(m_f) \circ f \ltsq{c_i}{c_o} i(m_g) \circ g$.
%
We make the analogous definition for error domains.

% We call such a square a \emph{perturbed square} to distinguish it from an ordinary square.
\end{definition}

Perturbed squares can be composed vertically and horizontally, as shown in
the following two lemmas. See the Appendix for proofs.

\begin{lemma}
  Suppose $f_1 \ltsqdelay{c_1}{c_2} g_1$ and $f_2 \ltsqdelay{c_2}{c_3} g_2$.
  Then $(f_2 \circ f_1) \ltsqdelay{c_1}{c_3} (g_2 \circ g_1)$.
\end{lemma}


\begin{lemma}
  If $f \ltsqdelay{c_i}{c_o} g$ and $g \ltsqdelay{c_i'}{c_o'} h$, then
  $f \ltsqdelay{c_i \odot c_i'}{c_o \odot c_o'} h$.
\end{lemma}

The important thing to note here is that the horizontal composition uses the
fact that the monoids are commutative, as well as the push-pull property of the
relations. For details, see the proof in the appendix.


%%%%%%%%%%%%%%%%%%%%%%%%%%%%%%%%%%%%%%%%%%%%%%%%%%%%%%%%%%%%%%%%%%%%%%%%%%%%%
\subsection{Horizontal Morphisms}

\subsubsection{Push-Pull}

As we have seen in the previous section, the push-pull property is necessary to
allow us to compose squares horizontally.  But as it turns out, our new
definition of vertical morphisms allows us to \emph{derive} push-pull for
relations that happen to be the restriction of an identity relation along a
morphism.

\begin{definition}
  Let $f : A \to A'$ be a perturbation-preserving morphism.
  The \textbf{(left-)restriction} of $r(A')$ along $f$, denoted $\restrl{r(A')}{f} : A \rel A'$, 
  is defined by $x \restrl{r(A')}{f} y := f\, x \mathrel{\ltdyn_{A'}} y$.

  Likewise let $g : A' \to A'$ be a perturbation-preserving morphism.
  The \textbf{(right-)restriction} of $r(A)$ along $g$, denoted $\restrr{r(A)}{g} : A \rel A'$,
  is defined by $x \restrr{r(A)}{g} y := x \mathrel{\ltdyn_{A}} g\, y$.
\end{definition}

\begin{lemma}\label{lem:restriction-push-pull}
  Let $f : A \to A'$ be a perturbation-preserving morphism. Then the
  relation $\restrl{r(A')}{f}$ has the push-pull property.

  Similarly, let $g : A' \to A$ be a perturbation-preserving morphism.
  Then the relation $\restrr{r(A)}{g}$ has the push-pull property.
\end{lemma}


\subsubsection{Representability}

Recall that the notion of \emph{quasi}-representability of a relation involves a
perturbation corresponding to the morphism opposite the cast. For each of the
\upl, \upr, \dnl, and \dnr rules there is a distinguished perturbation that
appears in the relevant square.
%
In the double category of casts we have just defined, the perturbations are
moved into the definition of the squares themselves. Thus, we can define
representability of a relation as the familiar notion of representability of a
horizontal morphism in a double category. In particular:

\begin{definition}
  Let $A$ and $A'$ be value objects, and let $c : A \rel A'$ be a
  predomain relation.  We say that $c$ is \emph{left-representable by}
  a perturbation-preserving morphism $e_c : A \to A'$ if there are
  squares $\upl : e_c \ltsqdelay{c}{r(A')} \id_{A'}$ and $\upr : \id_A
  \ltsqdelay{r(A)}{c} e_c$. We say that $c$ is \emph{left-representable}
  if it is left-representable by $e_c$ for some $e_c$.

  We say that $c$ is \emph{right-representable by} a
  perturbation-preserving morphism $p_c : A' \to A$ if there are
  squares $\dnr : \id_A \ltsqdelay{c}{r(A)} p_c$ and
  $\dnl : p_c \ltsqdelay{r(A')}{c} \id_{A'}$.

  We define the notion of representability for error domain relations
  analogously.
  
\end{definition}

An important consequence of this definition is that if $c$ is left (resp. right)
representable, then it is actually determined \emph{up to perturbation} by the
embedding (resp. projection) morphism.  More specifically, suppose $c$ is
left-representable by a morphism $e_c$. Recall the relation
$\restrl{r(A')}{e_c}$, i.e., the left restriction of the reflexive relation on
$A'$ along the morphism $e_c$. The fact that $c$ is left-representable by $e_c$
implies that $c$ is \emph{order-equivalent} to $\restrl{r(A')}{e_c}$, as shown
in the lemma below.

\begin{lemma}\label{lem:left-rep-implies-order-equiv}
  Let $c : A \rel A'$ be a predomain relation satisfying push-pull.
  If $c$ is left-representable by $e_c$, then there is a perturbed square 
  $\id \ltsqdelay{c}{\restrl{r(A')}{e_c}} \id$ and a perturbed square 
  $\id \ltsqdelay{\restrl{r(A')}{e_c}}{c} \id$.

  The analogous result is true for any error domain relation $d : B \rel B'$
  satisfying push-pull.
\end{lemma}

Similarly, right-representability by a morphism $p_c$ implies order equivalence
with the right restriction by $p_c$:

\begin{lemma}\label{lem:right-rep-implies-order-equiv} 
  Let $c : A \rel A'$ be a predomain relation satisfying push-pull.
  If $c$ is right-representable by $p_c$, then there is a perturbed square 
  $\id \ltsqdelay{c}{\restrr{r(A)}{p_c}} \id$ and a perturbed square 
  $\id \ltsqdelay{\restrr{r(A)}{p_c}}{c} \id$.

  The analogous result is true for any error domain relation $d : B \rel B'$
  satisfying push-pull.
\end{lemma}

The perturbations witnessing the order equivalence are hidden in the squares;
from the viewpoint of the double category, the two relations are order
equivalent.

\begin{remark}
  It is \emph{not} necessarily the case that if a relation $c : A \rel A'$ is
  left-representable by $e_c$, then it is actually \emph{equal} to the
  restriction $\restrl{r(A')}{e_c}$. Nor is a right-representable relation
  necessarily equal to a restriction.
  %
  For a simple example, consider the relation $F \iarr : FU(D \to FD) \rel FD$.
  This is quasi-right-representable in the old semantics, and hence it is
  right-representable in the new semantics where the perturbation is now part of
  the square.
  % But the projection morphism $p$ waits a step for the
  % later-function to become available. Thus,
  But the right-restriction of the reflexive relation $FU(D \to FD)$ along the
  projection morphism $p$ is not the same relation as $F \iarr$. Indeed, the
  restriction relates $x : FU(D \to FD)$ to $y : FD$ if and only if $x
  \ltdyn_{FU(D \to FD)} (p\, y)$. As the right-hand side inserts a $\theta$, the
  left-hand side must also be of the form $\theta$. But this means that the
  restriction relates $x$ of the form $\theta(\dots)$ to $y$ of the form
  $\eta(\dots)$. On the other hand, $F\iarr$ only relates $\theta$ to $\theta$.
  % the
  % former would need to relate a returned function on the left to a delayed function on the
  % right, whereas the latter requires both sides to be in lock-step.
\end{remark}


\subsubsection{Value and Computation Relations}

We now revisit the definition of value and computation relations. Recall that
the original definitions mentioned quasi-representability. We begin by restating
the definitions in terms of representability in the new category:

\begin{definition}\label{def:value-and-computation-relations}
  Let $A$ and $A'$ be value objects. A \textbf{value relation} $c : A \rel A'$
  is a relation between the underlying predomains such that:
  \begin{itemize}
  \item $c$ satisfies push-pull.
  \item $c$ is left-representable.
  \item $Fc$ is right-representable.
  \end{itemize}

  A \textbf{computation relation} $d : B \rel B'$ is defined analogously: an
  error domain relation satisfying push-pull, such that $d$ is
  right-representable and $Ud$ is left-representable.
\end{definition}

\begin{remark}
  The above definition is equivalent to the usual definition of embedding-projection pair.
  
  Specifically, suppose we are given an embedding $e : A \to A'$ and a projection $p
  : FA' \to FA$ satisfying $\id_{FA} \ltdyn p \circ Fe$ and $Fe \circ p \ltdyn
  \id_{FA'}$. Then we can recover a horizontal morphism in the above sense.
  %
  In particular, let $c$ be the relation $\restrl{r(A')}{e}$ and let $d$ be
  $\restrr{r(FA)}{p}$. By Lemma \ref{lem:restriction-push-pull},
  $c$ and $d$ satisfy the push-pull property.
  %
  It is also true that $c$ is left-representable (by $e$) and $d$ is
  right-representable (by $p$).

  We now show $Fc$ is order-equivalent to $d$ as follows:

% https://q.uiver.app/#q=WzAsOCxbMSwwLCJGQSJdLFsxLDEsIkZBJyJdLFsxLDIsIkZBIl0sWzIsMCwiRkEnIl0sWzIsMSwiRkEnIl0sWzIsMiwiRkEnIl0sWzAsMCwiRkEiXSxbMCwyLCJGQSJdLFswLDMsIkZjIiwwLHsic3R5bGUiOnsiYm9keSI6eyJuYW1lIjoiYmFycmVkIn0sImhlYWQiOnsibmFtZSI6Im5vbmUifX19XSxbNCwxLCJyKEZBJykiLDIseyJzdHlsZSI6eyJib2R5Ijp7Im5hbWUiOiJiYXJyZWQifSwiaGVhZCI6eyJuYW1lIjoibm9uZSJ9fX1dLFsyLDUsImQiLDIseyJzdHlsZSI6eyJib2R5Ijp7Im5hbWUiOiJiYXJyZWQifSwiaGVhZCI6eyJuYW1lIjoibm9uZSJ9fX1dLFswLDEsIkZlIiwyXSxbMSwyLCJwIiwyXSxbMyw0LCJcXGlkIl0sWzQsNSwiXFxpZCJdLFs3LDIsInIoRkEpIiwyLHsic3R5bGUiOnsiYm9keSI6eyJuYW1lIjoiYmFycmVkIn0sImhlYWQiOnsibmFtZSI6Im5vbmUifX19XSxbNiwwLCJyKEZBKSIsMCx7InN0eWxlIjp7ImJvZHkiOnsibmFtZSI6ImJhcnJlZCJ9LCJoZWFkIjp7Im5hbWUiOiJub25lIn19fV0sWzYsNywiXFxpZCIsMl0sWzExLDEzLCJGXFx1cGxfYyIsMSx7InNob3J0ZW4iOnsic291cmNlIjoyMCwidGFyZ2V0IjoyMH0sInN0eWxlIjp7ImJvZHkiOnsibmFtZSI6Im5vbmUifSwiaGVhZCI6eyJuYW1lIjoibm9uZSJ9fX1dLFsxMiwxNCwiXFxkbmxfZCIsMSx7InNob3J0ZW4iOnsic291cmNlIjoyMCwidGFyZ2V0IjoyMH0sInN0eWxlIjp7ImJvZHkiOnsibmFtZSI6Im5vbmUifSwiaGVhZCI6eyJuYW1lIjoibm9uZSJ9fX1dXQ==
\[\begin{tikzcd}[ampersand replacement=\&,sep=large]
	FA \& FA \& {FA'} \\
	\& {FA'} \& {FA'} \\
	FA \& FA \& {FA'}
	\arrow["{r(FA)}", "\shortmid"{marking}, no head, from=1-1, to=1-2]
	\arrow["\id"', from=1-1, to=3-1]
	\arrow["Fc", "\shortmid"{marking}, no head, from=1-2, to=1-3]
	\arrow[""{name=0, anchor=center, inner sep=0}, "Fe"', from=1-2, to=2-2]
	\arrow[""{name=1, anchor=center, inner sep=0}, "\id", from=1-3, to=2-3]
	\arrow[""{name=2, anchor=center, inner sep=0}, "p"', from=2-2, to=3-2]
	\arrow["{r(FA')}"', "\shortmid"{marking}, no head, from=2-3, to=2-2]
	\arrow[""{name=3, anchor=center, inner sep=0}, "\id", from=2-3, to=3-3]
	\arrow["{r(FA)}"', "\shortmid"{marking}, no head, from=3-1, to=3-2]
	\arrow["d"', "\shortmid"{marking}, no head, from=3-2, to=3-3]
	\arrow["{F\upl_c}"{description}, draw=none, from=0, to=1]
	\arrow["{\dnl_d}"{description}, draw=none, from=2, to=3]
\end{tikzcd}\]


Now since $d$ is right-representable, $Fc$ is also right-representable.
  
  
  % But we have no way of relating $Fc$ and $d$!

  % If, however, $e$ and $p$ satisfied the retraction property $p \circ Fe =
  % \id_{FA}$, then by vertically composing squares for representability of $Fc$
  % with those of $d$, we would obtain that $Fc$ and $d$ are order equivalent.
  % This would in turn imply that $Fc$ is right-representable, as is required by
  % the definition of horizontal morphism.

  % Unfortunately, in our semantics the retraction property is not true: it only
  % holds up to perturbation. We do not currently track this information in our
  % semantic model, leaving it instead to future exploration.
\end{remark}

% Just having an embedding and projection isn't enough to recover a horizontal
% morphism. Given such a pair (e, p), let c be the restriction of r(A') along e and let d be the restriction of r(FA) along p.
% We have that the restriction of r(A') along e is left-representable. But we have no way of relating 
%
% However, if the embedding and projection satisfied retraction, then we would
% be able to show that in fact the restriction of r(FA)

% However, by Lemmas \ref{lem:left-rep-implies-order-equiv} and
% \ref{lem:right-rep-implies-order-equiv} above, this definition is equivalent to
% the following:

% \begin{definition}
%   Let $A$ and $A'$ be value objects. A \textbf{value relation} $c : A \rel A'$
%   is given by a pair of an embedding morphism $e_c : A \to A'$ and a projection
%   morphism $p_c : FA' \to FA$.
% \end{definition}

% Suppose $c$ is a value relation according to the first definition. Then we can
% obtain a value relation according to the second definition: the fact that $c$ is
% left-representable means it has an embedding morphism $e_c$, and likewise
% there is a projection morphism $p_c$ for $Fc$.

% For the other direction, suppose we are given morphisms $e$ and $p$. We know by Lemma
% \ref{lem:restriction-push-pull} that the relations $\restrl{r(A')}{e}$ and
% $\restrr{r(FA)}{p}$ satisfy the push-pull property. Now consider the relation
% $c$ defined to be $\restrl{r(A')}{e}$. It is clear that $c$ is left-representable by $e$.


\subsection{Functorial Actions}

\subsection{Monadic Bind}

\subsection{Summary of the Model}

\textbf{Values}
\begin{itemize}
  \item Objects: Value objects $(A, M_A, i_A)$ where $A$ is a predomain, $M_A$
  is a commutative monoid, and $i_A : M_A \to \semptb{A}$ is a homomorphism of
  monoids. 
  \item Vertical morphisms: Morphisms of the underlying predomains that have an
  action on perturbations.
  \item Horizontal morphisms: Relations $c$ between the underlying predomains
  satisfying push-pull, such that $c$ is left-representable and $Fc$ is
  right-representable.
  \item Squares: Perturbed squares $f \ltsqdelay{c_i}{c_o} g$
\end{itemize}

\textbf{Computations}
\begin{itemize}
  \item Objects: Computation objects $(B, M_B, i_B)$
  \item Vertical morphisms: Morphisms of the underlying error domains that have an action on perturbations
  \item Horizontal morphisms: Relations $d$ between the underlying error domains
  satisfying push-pull, such that $d$ is right-representable and $Ud$ is
  left-representable
  \item Squares: Perturbed squares
\end{itemize}



%%%%%%%%%%%%%%%%%%%%%%%%%%%%%%%%%%%%%%%%%%%%%%%%%%%%%%%%%%%%%%%%%%%%%%%%%%%%%
\subsection{Revisiting the Graduality Proof}

Having described the double category in which the casts live, we now
return to the graduality proof.  We need to verify that the morphisms
of predomains/error domains used to define the cast morphisms do in
fact extend to perturbations.

We also need to check that the type constructors act on squares.


\subsubsection{The Model Construction}

\begin{enumerate}
  \item If $c$ and $c'$ are value relations, then there is a value relation $c \odot c'$. Likewise for computation relations.
  \item If $c$ is a value relation, then there is a computation relation $Fc$.
  \item If $d$ is a computation relation, then there is a value relation $Ud$.
  \item If $c_1$ and $c_2$ are value relations, then there is a value relation $c_1 \times c_2$
  \item If $c$ is a value relation and $d$ is a computation relation, then there is a computation relation $c \to d$.
\end{enumerate}


\appendix

\section{Omitted Examples and Proofs}

We begin with an example showing a concrete pair of type precision $c$ and $c'$
for which the projection for the composition $cc'$ inserts an additional
perturbation.

Let
%
$A_1 = (\dyntodyn) \ra \dyn$ and 
$A_2 = \dyntodyn$ and 
$A_3 = \dyn$,
%
and let
% 
$c = \iarr \ra r(\dyn) : A_1 \ltdyn A_2$ and
$c' = \iarr$.

First note that in the semantics, $F(cc')$ and $Fc Fc'$ are not equal as
relations, but they are quasi-equivalent since they are both
quasi-left-represented by the same embedding $Fe_{c'} \circ Fe_c$ (see
Lemma D.11 in the appendix of the POPL 2025 paper for details). 
%
In particular, there are syntactic perturbations $\delta^l \in M_{FA_1}$ and
$\delta^r \in M_{FA_3}$ and a square $\alpha : i(\delta^l) \ltsq{Fcc'}{FcFc'}
i(\delta^r)$.

Consider the $\dnr$ square for $F(cc')$, which says that $i(\delta_1)
\ltsq{F(cc')}{r(FA_1)} p_{F(cc')}$ for some perturbation $\delta_1 \in
M_{FA_1}$.
%
As in the proof of Lemma D.13, to define the $\dnr$ square for $F(cc')$, we
paste the square $\alpha$ on top of the $\dnr$ square for $FcFc'$.



As in the proof, we then observe that we must define the projection morphism
$p_{F(cc')}$ to be the composition of $i(\delta^r)$ with the projection
$p_{FcFc'}$.

The perturbations $\delta^l$ and $\delta^r$ are obtained as in the proof of
Lemma D.7 in the appendix; we can make that construction concrete since we know
what $d$ and $d'$ are. In particular, let us focus on $\delta^r$ since that is
the one used in the definition of the projection morphism for $F(cc')$.
%
Unfolding the proof of Lemma D.7, we see that the square $\alpha$ is obtained as
the following horizontal composition:

% https://q.uiver.app/#q=WzAsNixbMCwwLCJGQV8xIl0sWzEsMCwiRkFfMSJdLFsyLDAsIkZBXzMiXSxbMCwxLCJGQV8xIl0sWzEsMSwiRkFfMyJdLFsyLDEsIkZBXzMiXSxbMSw0LCJlIl0sWzAsMywiaShcXGRlbGxlX3tGY0ZjJ30pIiwyXSxbMiw1LCJpKFxcZGVscmVfe0YoY2MnKX0pIl0sWzAsMSwicihGQV8xKSIsMCx7InN0eWxlIjp7ImJvZHkiOnsibmFtZSI6ImJhcnJlZCJ9LCJoZWFkIjp7Im5hbWUiOiJub25lIn19fV0sWzEsMiwiRihjYycpIiwwLHsic3R5bGUiOnsiYm9keSI6eyJuYW1lIjoiYmFycmVkIn0sImhlYWQiOnsibmFtZSI6Im5vbmUifX19XSxbMyw0LCJGY0ZjJyIsMix7InN0eWxlIjp7ImJvZHkiOnsibmFtZSI6ImJhcnJlZCJ9LCJoZWFkIjp7Im5hbWUiOiJub25lIn19fV0sWzQsNSwicihGQV8zKSIsMix7InN0eWxlIjp7ImJvZHkiOnsibmFtZSI6ImJhcnJlZCJ9LCJoZWFkIjp7Im5hbWUiOiJub25lIn19fV1d
\[\begin{tikzcd}[ampersand replacement=\&]
	{FA_1} \& {FA_1} \& {FA_3} \\
	{FA_1} \& {FA_3} \& {FA_3}
	\arrow["{r(FA_1)}", "\shortmid"{marking}, no head, from=1-1, to=1-2]
	\arrow["{i(\delle_{FcFc'})}"', from=1-1, to=2-1]
	\arrow["{F(cc')}", "\shortmid"{marking}, no head, from=1-2, to=1-3]
	\arrow["e", from=1-2, to=2-2]
	\arrow["{i(\delre_{F(cc')})}", from=1-3, to=2-3]
	\arrow["{FcFc'}"', "\shortmid"{marking}, no head, from=2-1, to=2-2]
	\arrow["{r(FA_3)}"', "\shortmid"{marking}, no head, from=2-2, to=2-3]
\end{tikzcd}\]

Here, $e$ is the common embedding $Fe_{c'} \circ Fe_c$.

Thus, the perturbation we are interested in, $\delta^r$, is $\delre_{F(cc')}$.
Recalling the action of the functor $F$ on quasi-representable relations,
$\delre_{F(cc')}$ is simply the (injection into the coproduct of the)
perturbation $\delre_{cc'}$. Now recalling the definition of composition of
quasi-representable relations, we have $\delre_{cc'} = \delre_{c'} \cdot
\push_{c'}(\delre_c)$.

Now we can substitute the concrete values of $c$ and $c'$. We have
$c = U(\iarr \to r(FD)) : U(U(D \to FD) \to FD) \rel U(D \to FD)$ and 
$c' = \iarr : U(D \to FD) \rel D$.

The perturbation $\delre_{\iarr}$ is the identity since the corresponding
embedding doesn't delay, but the perturbation $\delre_c$ is nontrivial. It
corresponds to the fact that in the embedding for $c$, we project in the domain
according to $F\iarr$. \emph{This} perturbation is an actual delay in the case
that the dynamic value is a function.

Thus, we have isolated the place in the projection for $F(cc')$ where an
additional step may be taken in the relational semantics as compared to the
original term semantics.


\subsection{Vertical and Horizontal Composition of Squares}

%%%%%%%%%%%%%%%%%%%%%%%%%%%%%%%%%%%
% Vertical composition of squares
%%%%%%%%%%%%%%%%%%%%%%%%%%%%%%%%%%%
\begin{lemma}
  Suppose $f_1 \ltsqdelay{c_1}{c_2} g_1$ and $f_2 \ltsqdelay{c_2}{c_3} g_2$.
  Then $(f_2 \circ f_1) \ltsqdelay{c_1}{c_3} (g_2 \circ g_1)$.
\end{lemma}

\begin{proof}
  Let $m_{f_1}$ and $m_{g_1}$ be the perturbations in the first
  square, and let $m_{f_2}$ and $m_{g_2}$ be the perturbations in the
  second square. We compose the two given squares vertically to obtain

\begin{comment}
% https://q.uiver.app/#q=WzAsNixbMCwwLCJBXzEiXSxbMCwxLCJBXzIiXSxbMCwyLCJBXzIiXSxbMiwwLCJBXzEnIl0sWzIsMSwiQV8yJyJdLFsyLDIsIkFfMiciXSxbMCwxLCJmXzEiLDJdLFsxLDIsImlfe0FfMn0obV97Zl8xfSkiLDJdLFszLDQsImdfMSJdLFs0LDUsImlfe0FfMid9KG1fe2dfMX0pIl0sWzAsMywiY18xIiwwLHsic3R5bGUiOnsiYm9keSI6eyJuYW1lIjoiYmFycmVkIn0sImhlYWQiOnsibmFtZSI6Im5vbmUifX19XSxbMiw1LCJjXzIiLDIseyJzdHlsZSI6eyJib2R5Ijp7Im5hbWUiOiJiYXJyZWQifSwiaGVhZCI6eyJuYW1lIjoibm9uZSJ9fX1dXQ==
\[\begin{tikzcd}[ampersand replacement=\&]
        {A_1} \&\& {A_1'} \\
        {A_2} \&\& {A_2'} \\
        {A_2} \&\& {A_2'}
        \arrow["{c_1}", "\shortmid"{marking}, no head, from=1-1, to=1-3]
        \arrow["{f_1}"', from=1-1, to=2-1]
        \arrow["{g_1}", from=1-3, to=2-3]
        \arrow["{i_{A_2}(m_{f_1})}"', from=2-1, to=3-1]
        \arrow["{i_{A_2'}(m_{g_1})}", from=2-3, to=3-3]
        \arrow["{c_2}"', "\shortmid"{marking}, no head, from=3-1, to=3-3]
\end{tikzcd}\]

and

% https://q.uiver.app/#q=WzAsNixbMCwwLCJBXzIiXSxbMCwxLCJBXzMiXSxbMCwyLCJBXzMiXSxbMiwwLCJBXzInIl0sWzIsMSwiQV8zJyJdLFsyLDIsIkFfMyciXSxbMCwxLCJmXzIiLDJdLFsxLDIsImlfe0FfM30obV97Zl8yfSkiLDJdLFszLDQsImdfMiJdLFs0LDUsImlfe0FfMyd9KG1fe2dfMn0pIl0sWzAsMywiY18yIiwwLHsic3R5bGUiOnsiYm9keSI6eyJuYW1lIjoiYmFycmVkIn0sImhlYWQiOnsibmFtZSI6Im5vbmUifX19XSxbMiw1LCJjXzMiLDIseyJzdHlsZSI6eyJib2R5Ijp7Im5hbWUiOiJiYXJyZWQifSwiaGVhZCI6eyJuYW1lIjoibm9uZSJ9fX1dXQ==
\[\begin{tikzcd}[ampersand replacement=\&]
	{A_2} \&\& {A_2'} \\
	{A_3} \&\& {A_3'} \\
	{A_3} \&\& {A_3'}
	\arrow["{c_2}", "\shortmid"{marking}, no head, from=1-1, to=1-3]
	\arrow["{f_2}"', from=1-1, to=2-1]
	\arrow["{g_2}", from=1-3, to=2-3]
	\arrow["{i_{A_3}(m_{f_2})}"', from=2-1, to=3-1]
	\arrow["{i_{A_3'}(m_{g_2})}", from=2-3, to=3-3]
	\arrow["{c_3}"', "\shortmid"{marking}, no head, from=3-1, to=3-3]
\end{tikzcd}\]
\end{comment}

% https://q.uiver.app/#q=WzAsMTAsWzAsMCwiQV8xIl0sWzAsMSwiQV8yIl0sWzAsMiwiQV8yIl0sWzIsMCwiQV8xJyJdLFsyLDEsIkFfMiciXSxbMiwyLCJBXzInIl0sWzAsMywiQV8zIl0sWzAsNCwiQV8zIl0sWzIsMywiQV8zJyJdLFsyLDQsIkFfMyciXSxbMCwxLCJmXzEiLDJdLFsxLDIsImlfe0FfMn0obV97Zl8xfSkiLDJdLFszLDQsImdfMSJdLFs0LDUsImlfe0FfMid9KG1fe2dfMX0pIl0sWzAsMywiY18xIiwwLHsic3R5bGUiOnsiYm9keSI6eyJuYW1lIjoiYmFycmVkIn0sImhlYWQiOnsibmFtZSI6Im5vbmUifX19XSxbMiw1LCJjXzIiLDIseyJzdHlsZSI6eyJib2R5Ijp7Im5hbWUiOiJiYXJyZWQifSwiaGVhZCI6eyJuYW1lIjoibm9uZSJ9fX1dLFsyLDYsImZfMiIsMl0sWzUsOCwiZ18yIl0sWzYsNywiaV97QV8zfShtX3tmXzJ9KSIsMl0sWzgsOSwiaV97QV8zJ30obV97Z18yfSkiXSxbNyw5LCJjXzMiLDIseyJzdHlsZSI6eyJib2R5Ijp7Im5hbWUiOiJiYXJyZWQifSwiaGVhZCI6eyJuYW1lIjoibm9uZSJ9fX1dXQ==
\[\begin{tikzcd}[ampersand replacement=\&]
	{A_1} \&\& {A_1'} \\
	{A_2} \&\& {A_2'} \\
	{A_2} \&\& {A_2'} \\
	{A_3} \&\& {A_3'} \\
	{A_3} \&\& {A_3'}
	\arrow["{c_1}", "\shortmid"{marking}, no head, from=1-1, to=1-3]
	\arrow["{f_1}"', from=1-1, to=2-1]
	\arrow["{g_1}", from=1-3, to=2-3]
	\arrow["{i_{A_2}(m_{f_1})}"', from=2-1, to=3-1]
	\arrow["{i_{A_2'}(m_{g_1})}", from=2-3, to=3-3]
	\arrow["{c_2}"', "\shortmid"{marking}, no head, from=3-1, to=3-3]
	\arrow["{f_2}"', from=3-1, to=4-1]
	\arrow["{g_2}", from=3-3, to=4-3]
	\arrow["{i_{A_3}(m_{f_2})}"', from=4-1, to=5-1]
	\arrow["{i_{A_3'}(m_{g_2})}", from=4-3, to=5-3]
	\arrow["{c_3}"', "\shortmid"{marking}, no head, from=5-1, to=5-3]
\end{tikzcd}\]

Now we commute the inner perturbation with the inner function on both
sides. That is, we have
%
$f_2 \circ i_{A_2}(m_{f_1}) = i_{A_3}(f_2^\to(m_{f_1})) \circ f_2$,
%
and likewise
%
$g_2 \circ i_{A_2'}(m_{g_1}) = i_{A_3'}(g_2^\to(m_{g_1})) \circ g_2$,
%
where $f_2^\to : M_{A_2} \to M_{A_3}$ is the homomorphism of monoids
corresponding to $f_2$ and similarly for $g_2^\to$.
Doing so, we see that the above square is equivalent to

% https://q.uiver.app/#q=WzAsMTAsWzAsMCwiQV8xIl0sWzAsMSwiQV8yIl0sWzAsMiwiQV8yIl0sWzIsMCwiQV8xJyJdLFsyLDEsIkFfMiciXSxbMiwyLCJBXzInIl0sWzAsMywiQV8zIl0sWzAsNCwiQV8zIl0sWzIsMywiQV8zJyJdLFsyLDQsIkFfMyciXSxbMCwxLCJmXzEiLDJdLFsxLDIsImZfMiIsMl0sWzMsNCwiZ18xIl0sWzQsNSwiZ18yIl0sWzAsMywiY18xIiwwLHsic3R5bGUiOnsiYm9keSI6eyJuYW1lIjoiYmFycmVkIn0sImhlYWQiOnsibmFtZSI6Im5vbmUifX19XSxbMiw2LCJpX3tBXzN9KGZfMl5cXHRvKG1fe2ZfMX0pKSIsMl0sWzUsOCwiaV97QV8zJ30oZ18yXlxcdG8obV97Z18xfSkpIl0sWzYsNywiaV97QV8zfShtX3tmXzJ9KSIsMl0sWzgsOSwiaV97QV8zJ30obV97Z18yfSkiXSxbNyw5LCJjXzMiLDIseyJzdHlsZSI6eyJib2R5Ijp7Im5hbWUiOiJiYXJyZWQifSwiaGVhZCI6eyJuYW1lIjoibm9uZSJ9fX1dXQ==
\[\begin{tikzcd}[ampersand replacement=\&]
	{A_1} \&\& {A_1'} \\
	{A_2} \&\& {A_2'} \\
	{A_2} \&\& {A_2'} \\
	{A_3} \&\& {A_3'} \\
	{A_3} \&\& {A_3'}
	\arrow["{c_1}", "\shortmid"{marking}, no head, from=1-1, to=1-3]
	\arrow["{f_1}"', from=1-1, to=2-1]
	\arrow["{g_1}", from=1-3, to=2-3]
	\arrow["{f_2}"', from=2-1, to=3-1]
	\arrow["{g_2}", from=2-3, to=3-3]
	\arrow["{i_{A_3}(f_2^\to(m_{f_1}))}"', from=3-1, to=4-1]
	\arrow["{i_{A_3'}(g_2^\to(m_{g_1}))}", from=3-3, to=4-3]
	\arrow["{i_{A_3}(m_{f_2})}"', from=4-1, to=5-1]
	\arrow["{i_{A_3'}(m_{g_2})}", from=4-3, to=5-3]
	\arrow["{c_3}"', "\shortmid"{marking}, no head, from=5-1, to=5-3]
\end{tikzcd}\]

Thus, we can take the perturbation corresponding to $f_2 \circ f_1$ to
be $m_{f_2} \cdot f_2^\to(m_{f_1})$ and the pertubation corresponding
to $g_2 \circ g_1$ to be $m_{g_2} \cdot g_2^\to(m_{g_1})$.

\end{proof}

%%%%%%%%%%%%%%%%%%%%%%%%%%%%%%%%%%%%%
% Horizontal composition of squares
%%%%%%%%%%%%%%%%%%%%%%%%%%%%%%%%%%%%%
\begin{lemma}
  If $f \ltsqdelay{c_i}{c_o} g$ and $g \ltsqdelay{c_i'}{c_o'} h$, then
  $f \ltsqdelay{c_i \odot c_i'}{c_o \odot c_o'} h$.
\end{lemma}
\begin{proof}

\end{proof}


\subsection{Lemmas about Push-Pull, Representability}

%%%%%%%%%%%%%%%%%%%%%%%%%%%%%%%%%%%%%%%%%%%%%%%%%%%%%%%%%%%%%%%%%%%
% Deriving push-pull for the restriction of a reflexive relation
%%%%%%%%%%%%%%%%%%%%%%%%%%%%%%%%%%%%%%%%%%%%%%%%%%%%%%%%%%%%%%%%%%%

% \begin{lemma}
%   Let $f : A \to A'$ be a perturbation-preserving morphism. Then the
%   relation $\restr^l{\ltdyn_{A'}}{f} : A \rel A'$, defined by $x \restr^l{\ltdyn_{A'}}{f}
%   y := f x \mathrel{\ltdyn_{A'}} y$, has the push-pull property.
% \end{lemma}

\begin{lemma}
  Let $f : A \to A'$ be a perturbation-preserving morphism. Then the
  relation $\restrl{r(A')}{f}$ has the push-pull property.

  Similarly, let $g : A' \to A$ be a perturbation-preserving morphism.
  Then the relation $\restrr{r(A)}{g}$ has the push-pull property.
\end{lemma}
\begin{proof}
  First we show how to derive push. The push homomorphism is simply
  given by $h^\to$; it remains to show the existence of the following
  square:

% https://q.uiver.app/#q=WzAsNCxbMCwwLCJBIl0sWzAsMSwiQSJdLFsxLDAsIkEnIl0sWzEsMSwiQSciXSxbMCwyLCJcXHJlc3RybHtyKEEnKX17Zn0iLDAseyJzdHlsZSI6eyJib2R5Ijp7Im5hbWUiOiJiYXJyZWQifSwiaGVhZCI6eyJuYW1lIjoibm9uZSJ9fX1dLFsxLDMsIlxccmVzdHJse3IoQScpfXtmfSIsMCx7InN0eWxlIjp7ImJvZHkiOnsibmFtZSI6ImJhcnJlZCJ9LCJoZWFkIjp7Im5hbWUiOiJub25lIn19fV0sWzAsMSwiaV9BKG0pIiwyXSxbMiwzLCJpX3tBJ30oaF57XFx0b30obSkpIl1d
\[\begin{tikzcd}[ampersand replacement=\&]
	A \& {A'} \\
	A \& {A'}
	\arrow["{\restrl{r(A')}{f}}", "\shortmid"{marking}, no head, from=1-1, to=1-2]
	\arrow["{i_A(m)}"', from=1-1, to=2-1]
	\arrow["{i_{A'}(h^{\to}(m))}", from=1-2, to=2-2]
	\arrow["{\restrl{r(A')}{f}}", "\shortmid"{marking}, no head, from=2-1, to=2-2]
\end{tikzcd}\]

By definition, it will suffice to construct a square

% https://q.uiver.app/#q=WzAsNixbMCwwLCJBIl0sWzAsMSwiQSJdLFsxLDAsIkEnIl0sWzAsMiwiQSciXSxbMSwyLCJBJyJdLFsxLDEsIkEnIl0sWzAsMiwiXFxyZXN0cmx7cihBJyl9e2Z9IiwwLHsic3R5bGUiOnsiYm9keSI6eyJuYW1lIjoiYmFycmVkIn0sImhlYWQiOnsibmFtZSI6Im5vbmUifX19XSxbMCwxLCJpX0EobSkiLDJdLFszLDQsInIoQScpIiwwLHsic3R5bGUiOnsiYm9keSI6eyJuYW1lIjoiYmFycmVkIn0sImhlYWQiOnsibmFtZSI6Im5vbmUifX19XSxbMSwzLCJmIiwyXSxbNSw0LCJcXGlkIl0sWzIsNSwiaV97QSd9KGhee1xcdG99KG0pKSJdXQ==
\[\begin{tikzcd}[ampersand replacement=\&]
	A \& {A'} \\
	A \& {A'} \\
	{A'} \& {A'}
	\arrow["{\restrl{r(A')}{f}}", "\shortmid"{marking}, no head, from=1-1, to=1-2]
	\arrow["{i_A(m)}"', from=1-1, to=2-1]
	\arrow["{i_{A'}(h^{\to}(m))}", from=1-2, to=2-2]
	\arrow["f"', from=2-1, to=3-1]
	\arrow["\id", from=2-2, to=3-2]
	\arrow["{r(A')}", "\shortmid"{marking}, no head, from=3-1, to=3-2]
\end{tikzcd}\]

Now since $f$ preserves perturbations, the left-hand side of the
square is equal to $i_{A'}(h^\to(m)) \circ f$, so the above is the same
as

% https://q.uiver.app/#q=WzAsNixbMCwwLCJBIl0sWzAsMSwiQSciXSxbMSwwLCJBJyJdLFswLDIsIkEnIl0sWzEsMiwiQSciXSxbMSwxLCJBJyJdLFswLDIsIlxccmVzdHJse3IoQScpfXtmfSIsMCx7InN0eWxlIjp7ImJvZHkiOnsibmFtZSI6ImJhcnJlZCJ9LCJoZWFkIjp7Im5hbWUiOiJub25lIn19fV0sWzAsMSwiZiIsMl0sWzMsNCwicihBJykiLDAseyJzdHlsZSI6eyJib2R5Ijp7Im5hbWUiOiJiYXJyZWQifSwiaGVhZCI6eyJuYW1lIjoibm9uZSJ9fX1dLFsxLDMsImlfe0EnfShoXntcXHRvfShtKSkiLDJdLFs1LDQsImlfe0EnfShoXntcXHRvfShtKSkiXSxbMiw1LCJcXGlkIl0sWzEsNSwicihBJykiLDAseyJzdHlsZSI6eyJib2R5Ijp7Im5hbWUiOiJiYXJyZWQifSwiaGVhZCI6eyJuYW1lIjoibm9uZSJ9fX1dXQ==
\[\begin{tikzcd}[ampersand replacement=\&]
	A \& {A'} \\
	{A'} \& {A'} \\
	{A'} \& {A'}
	\arrow["{\restrl{r(A')}{f}}", "\shortmid"{marking}, no head, from=1-1, to=1-2]
	\arrow["f"', from=1-1, to=2-1]
	\arrow["\id", from=1-2, to=2-2]
	\arrow["{r(A')}", "\shortmid"{marking}, no head, from=2-1, to=2-2]
	\arrow["{i_{A'}(h^{\to}(m))}"', from=2-1, to=3-1]
	\arrow["{i_{A'}(h^{\to}(m))}", from=2-2, to=3-2]
	\arrow["{r(A')}", "\shortmid"{marking}, no head, from=3-1, to=3-2]
\end{tikzcd}\]

The top square is the definition of the restriction of a relation
along a morphism, and the second square is a horizontal identity
square. Composing these vertically, we are finished.

The pull property is derived in an analogous manner.
  
\end{proof}


%%%%%%%%%%%%%%%%%%%%%%%%%%%%%%%%%%%%%%%%%%%%%%%%%%%%%%%%%%%%%%%%%%
% Push-pull + representability
%   ==> order equivalent to restriction of a reflexive relation
%%%%%%%%%%%%%%%%%%%%%%%%%%%%%%%%%%%%%%%%%%%%%%%%%%%%%%%%%%%%%%%%%%
\begin{lemma}
  Let $c : A \rel A'$ be a predomain relation satisfying push-pull.
  If $c$ is left-representable by $e_c$, then there is a perturbed square 
  $\id \ltsqdelay{c}{\restrl{r(A')}{e_c}} \id$ and a perturbed square 
  $\id \ltsqdelay{\restrl{r(A')}{e_c}}{c} \id$.

  The analogous result is true for any error domain relation $d : B \rel B'$
  satisfying push-pull.
\end{lemma}
\begin{proof}
  We first show that $\restrl{r(A')}{e_c}$ is left-representable by $e_c$. The
  square $\upl : e_c \ltsq{\restrl{r(A')}{e_c}}{r(A')} \id$ exists by definition
  of the restriction. The square $\upr : \id \ltsq{r(A)}{\restrl{r(A')}{e_c}}
  e_c$ is equivalent to the square $e_c \ltsq{r(A)}{r(A')} e_c$ which is just an
  identity square.

  Then since $c$ is also left-representable by $e_c$, we can combine the squares
  horizontally to obtain the desired squares:

  % https://q.uiver.app/#q=WzAsNixbMCwwLCJBIl0sWzEsMCwiQSJdLFsyLDAsIkEnIl0sWzAsMSwiQSJdLFsxLDEsIkEnIl0sWzIsMSwiQSciXSxbMCwzLCJcXGlkIiwyXSxbMSw0LCJlX2MiXSxbMiw1LCJcXGlkIl0sWzAsMSwicihBKSIsMCx7InN0eWxlIjp7ImJvZHkiOnsibmFtZSI6ImJhcnJlZCJ9LCJoZWFkIjp7Im5hbWUiOiJub25lIn19fV0sWzEsMiwiYyIsMCx7InN0eWxlIjp7ImJvZHkiOnsibmFtZSI6ImJhcnJlZCJ9LCJoZWFkIjp7Im5hbWUiOiJub25lIn19fV0sWzMsNCwiXFxyZXN0cmx7cihBJyl9e2VfY30iLDIseyJzdHlsZSI6eyJib2R5Ijp7Im5hbWUiOiJiYXJyZWQifSwiaGVhZCI6eyJuYW1lIjoibm9uZSJ9fX1dLFs0LDUsInIoQScpIiwyLHsic3R5bGUiOnsiYm9keSI6eyJuYW1lIjoiYmFycmVkIn0sImhlYWQiOnsibmFtZSI6Im5vbmUifX19XSxbNyw4LCIoMikiLDMseyJzaG9ydGVuIjp7InNvdXJjZSI6MjAsInRhcmdldCI6MjB9LCJzdHlsZSI6eyJib2R5Ijp7Im5hbWUiOiJub25lIn0sImhlYWQiOnsibmFtZSI6Im5vbmUifX19XSxbNiw3LCIoMSkiLDMseyJzaG9ydGVuIjp7InNvdXJjZSI6MjAsInRhcmdldCI6MjB9LCJzdHlsZSI6eyJib2R5Ijp7Im5hbWUiOiJub25lIn0sImhlYWQiOnsibmFtZSI6Im5vbmUifX19XV0=
  \[\begin{tikzcd}[ampersand replacement=\&,column sep=3.15em]
    A \& A \& {A'} \\
    A \& {A'} \& {A'}
    \arrow["{r(A)}", "\shortmid"{marking}, no head, from=1-1, to=1-2]
    \arrow[""{name=0, anchor=center, inner sep=0}, "\id"', from=1-1, to=2-1]
    \arrow["c", "\shortmid"{marking}, no head, from=1-2, to=1-3]
    \arrow[""{name=1, anchor=center, inner sep=0}, "{e_c}", from=1-2, to=2-2]
    \arrow[""{name=2, anchor=center, inner sep=0}, "\id", from=1-3, to=2-3]
    \arrow["{\restrl{r(A')}{e_c}}"', "\shortmid"{marking}, no head, from=2-1, to=2-2]
    \arrow["{r(A')}"', "\shortmid"{marking}, no head, from=2-2, to=2-3]
    \arrow["{(1)}"{marking, allow upside down}, draw=none, from=0, to=1]
    \arrow["{(2)}"{marking, allow upside down}, draw=none, from=1, to=2]
  \end{tikzcd}\]

  Here, (1) is the $\upr$ square for $\restrl{r(A')}{e_c}$ while (2) is the $\upl$ square for $c$.
  The other square is constructed symmetrically.


  The proof for error domain relations is analogous.

  % Need to paste the relevant squares horizontally along the equal
  % embedding, but we cannot compose squares horizontally without the
  % relations satisfying push-pull.
  
\end{proof}

\begin{lemma}
  Let $c : A \rel A'$ be a predomain relation satisfying push-pull.
  If $c$ is right-representable by $p_c$, then there is a perturbed square 
  $\id \ltsqdelay{c}{\restrr{r(A)}{p_c}} \id$ and a perturbed square 
  $\id \ltsqdelay{\restrr{r(A)}{p_c}}{c} \id$.

  The analogous result is true for any error domain relation $d : B \rel B'$
  satisfying push-pull.
\end{lemma}
\begin{proof}
  Analogous to the proof of the previous lemma.
\end{proof}

\end{document}

